\documentclass[12pt, a4paper]{article}

\usepackage{amsmath}
\usepackage[utf8]{inputenc}
\usepackage[spanish]{babel} %Paquete de idioma
\usepackage[hidelinks]{hyperref}
\usepackage{graphicx}
\usepackage{float}
\usepackage{eso-pic}
\usepackage{lipsum}
\usepackage{transparent}
\usepackage{parskip}
%\usepackage[backend=biber, style=apa]{biblatex}

\graphicspath{{imagenes/}}

%%%%%%%%% ESTO PARA LA MARCA DE AGUA %%%%%%%%%%%%%%%%%%%%%%%%%%%%%%%


\AddToShipoutPicture{ 
    \put(410,380){
        \parbox[b][\paperheight]{\paperwidth}{%
            \vfill
            \
            {
            \transparent{1}
            \includegraphics[scale=0.5]{logo-ua.png}
            \vfill
            }
        }
    }
}
%%%%%%%%%%%%%%%%%%%%%%%%%%%%%%%%%%%%%%%%%%%%%%%%%%%%%%%%%%

\title{Memoria Práctica 2 \\


\large Comunicaciones
}
\author{
Leopoldo Cadavid Piñero
}
\date{Febrero 2022}







\begin{document}

\maketitle
\newpage
\tableofcontents
\newpage
\section{Introducción}
      

\section{Sesión 1}

\section{Sesión 2}
 En esta sesión se plantean cómo objetivos:
 \begin{itemize}
     \item Terminar con las pruebas del programa de la sesión 1, Ubidots
     \item Iniciar el porceso de automatización de una sencilla instalación, usando protocolos de 
        comunicación vistos en la teoría.
    
    
    \end{itemize}

\subsection{Iniciación de control automatico}

Pasos:
\begin{itemize}
    \item Conectarse a la red del routeer de la instalación que es la subred de control
    \item meterse en la  interfaz de shelly
    \item Desde aquí podremos conectarnos a los dispositivos del sistema y probar como podemos dar instrucciones y recibir info
    \item \textbf{Aconsejable:} decirle al shelly a que red se tiene que conectar. Decimos a shelly que se conecte a la red que nosotros queremos
    (esto se graba en su memoria).
    \item Para nuestro caso: \textbf{Nombre de la red:} Cudy-081c, \textbf{IP:} 192.168.10.100 (rango x.x.10.0/255), 
    \textbf{Gateway (opcional):}. 
    \item Si hemos introducido todo bien, nuestro shell estará conectado en la red del router.

\end{itemize}











\end{document}